\documentclass[uplatex,a4paper,12pt]{jsarticle}
\usepackage{amsmath,amssymb}
\usepackage{enumitem}

\begin{document}

\title{ニューラルネットワークの最適化・正則化 確認テスト 解答}
\author{}
\date{}
\maketitle

\section*{配点}
全10問，各10点，合計100点．60点以上を合格とする．

\section*{解答}

\begin{enumerate}

\item
\[
\boxed{(b)}
\]
SGDは，一部のサンプル（ミニバッチ）を用いて更新を行う手法である．

\item
\[
\boxed{\text{エポック}}
\]

\item
\[
\boxed{(a)}
\]
モメンタム法では，前回の更新量を利用して学習を安定化・高速化する．

\item
\[
\boxed{(a)}
\]
AdamはRMSPropを発展させた最適化手法であり，深層学習で広く利用される．

\item
\[
\boxed{(b)}
\]
Xavier初期化は，sigmoidやtanhを想定した初期化方法である．

\item
\[
\boxed{(c)}
\]
He初期化はReLUを用いる場合を想定した初期化である．

\item
\[
\boxed{(a)}
\]
バッチ正規化は内部共変量シフトを抑える目的で用いられる．

\item
\[
\boxed{\text{重み減衰}}
\]
または
\[
\boxed{\text{Weight Decay}}
\]

\item
\[
\boxed{(a)}
\]
Dropoutは，一部ユニットをランダムに無効化することで過学習を抑える．

\item
\[
\boxed{(a)}
\]
転移学習では，他タスクで学習した知識やパラメータを利用する．

\end{enumerate}

\end{document}